Dues càrregues elèctriques puntuals de $-5,0\ \mu\mathrm{C}$ i $7,0\ \mu\mathrm{C}$ estan separades 10~cm l'una de l'altra.

\begin{apartat}{1,25}
Calculeu el camp elèctric (mòdul, direcció i sentit) en un punt a 3,0~cm de la càrrega negativa i a 7,0~cm de la càrrega positiva. Aquest punt pertany a la línia que uneix les dues càrregues.
\end{apartat}

\begin{apartat}{1,25}
Calculeu en quin punt de la línia que uneix les càrregues el potencial elèctric és nul.
\end{apartat}

\blocetiquetat{Dada:}{%
  $k = \dfrac{1}{4\pi\varepsilon_0} = 8,99\times 10^{9}\ \mathrm{N\,m^2\,C^{-2}}$.}
